\documentclass[a4paper,oneside,12pt]{amsart}

\usepackage[spanish]{babel}
\usepackage[utf8]{inputenc}
%\usepackage[latin1]{inputenc}
\usepackage{latexsym}
\usepackage{multicol}
\usepackage{enumerate}
\usepackage[heightrounded]{geometry}
\geometry{top=2cm,bottom=2cm,left=2cm,right=2cm}
\usepackage{graphicx}

\DeclareMathOperator{\cond}{cond}
\DeclareMathOperator{\Id}{Id}
\DeclareMathOperator{\re}{Re}
%\DeclareMathOperator{\im}{Im}
\newcommand{\I}{\mathbb{I}}
\newcommand{\R}{\mathbb{R}}
\newcommand{\Q}{\mathbb{Q}}
\newcommand{\C}{\mathbb{C}}
\newcommand{\Z}{\mathbb{Z}}
\newcommand{\N}{\mathbb{N}}
\DeclareMathOperator{\im}{Im}

\newcommand{\nombremateria}
{Investigaci\'on Operativa}
\newcommand{\nombrecuatrimestre}
    {Segundo cuatrimestre de 2020}
\newcommand{\numeroparcial}
    { Recuperatorio del Segundo Parcial }
\newcommand{\fecha}
    {23 de Diciembre de 2020}
\newcommand{\numerotema}
    {\bf 1}

%\oddsidemargin -1cm \evensidemargin -1cm \textwidth 17.5cm
%topmargin 1.2cm \textheight 25cm
%\setlength{\textheight}{25cm}

\def \ds{\displaystyle}
\newtheorem{quest}{Ejercicio}
\thispagestyle{empty}

\begin{document}

\hfill \hfill
\begin{tabular}[c]{|c|c|c|c|c|}
\hline 
1 & 2 & 3 & 4 & Calificaci\'on \\ \hline \hline
\hspace{1.2 cm} & \hspace{1.2 cm} & \hspace{1.2 cm} & \hspace{1.2 cm} &\hspace{1.5 cm} \\
\hspace{1.2 cm} & \hspace{1.2 cm} & \hspace{1.2 cm} & \hspace{1.2 cm} &\hspace{1.5 cm} \\ \hline

\end{tabular}

\vspace{-1cm}

\noindent Nombre y apellido:

\vspace{0.2cm}

\noindent Número de libreta:

\noindent \hrulefill 

\smallskip
\renewcommand{\baselinestretch}{1.1}

\begin{center}
	\large \textbf{\nombremateria} 
	
	\numeroparcial -- \fecha
\end{center}

\vspace{.6cm}
\renewcommand{\theenumi}{\textbf{\arabic{enumi}}}

%%% EMPIEZA EL EJERCICIO 1 %%%

\begin{quest} 
Resolver el siguiente problema de Programación Lineal utilizando el Método \mbox{Big-M}. Explicitar dónde se realiza el óptimo y el valor que toma la función objetivo en el mismo.
    \[\begin{array}{rrrrrrrr}
        \max & 2x_1 & + & x_2 & - & x_3  \\
        \text{s.a:} & 2x_1 & + & 3x_2 & + & x_3 & \leq & 9 \\
          &  &  & 2x_2 & - & x_3 & \geq & 4 \\
          & x_1  &  &  & - & x_3 & = & 6 \\
         & \multicolumn{7}{l}{x_3\leq 0 \;\;\;x_1,x_2\geq 0}
    \end{array}
    \]

\end{quest}

%%% TERMINA EL EJERCICIO 1 %%%

%%% EMPIEZA EL EJERCICIO 2 %%%

\begin{quest} 
Los siguientes diccionarios corresponden a una iteración de SIMPLEX de diferentes problemas lineales con objetivo de maximizar.

\begin{multicols}{2}
    \begin{itemize}
        \item[1)]
            \[
            \begin{array}{rrrrrrr}
                x_3&=&&&2x_2&+&x_5\\
                x_1&=&2&+&x_2&-&x_5\\
                x_4&=&1&-&x_2&-&2x_5\\ \hline
                z&=&3&-&x_2&-&\frac{1}{2}x_5
            \end{array}
            \]
        \item[3)]
            \[
            \begin{array}{rrrrrrr}
                x_2&=&\frac{1}{2}&+&x_1&-&2x_3\\
                x_5&=&5&-&x_1&+&x_3\\
                x_4&=&1&+&x_1&&\\ \hline
                z&=&11&&&-&\frac{3}{2}x_3
            \end{array}
            \]
        \item[2)]
            \[
            \begin{array}{rrrrrrr}
                x_1&=&2&&&-&x_4\\
                x_3&=&\frac{2}{3}&+&\frac{1}{2}x_2&+&\frac{1}{4}x_4\\
                x_5&=&3&+&2x_2&-&\frac{1}{2}x_4\\ \hline
                z&=&-4&+&\frac{1}{3}x_2&-&x_4
            \end{array}
            \]
        \item[4)]
            \[
            \begin{array}{rrrrrrr}
                x_5&=&\frac{3}{2}&+&x_2&+&\frac{1}{2}x_3\\
                x_4&=&1&-&x_2&-&x_3\\
                x_1&=&\frac{1}{10}&+&2x_2&-&\frac{4}{3}x_3\\ \hline
                z&=&&-&x_2&-&\frac{2}{3}x_3
            \end{array}
            \]
    \end{itemize}
\end{multicols}
Indicar y justificar a qué situación corresponde cada uno de ellos:
\begin{enumerate}[a)]
    \item Problema lineal con infinitas soluciones.
    \item Problema lineal con solución única no degenerada.
    \item Problema lineal no acotado.
    \item Problema lineal con solución única degenerada.
\end{enumerate}
\end{quest}

%%% TERMINA EL EJERCICIO 2 %%%

%%% EMPIEZA EL EJERCICIO 3 %%%


\begin{quest} 
Dado el siguiente problema de programación lineal:

\[\begin{array}{rrrrrrrrrrrrrr}
    \max & 4x_1 & + & x_2 & - & 2x_3 & + & x_4 & - & x_5 & + & 6x_6 & & \\
    \text{s.a.} & -x_1 & + & 3x_2 & - & x_3 & - & x_4 & + & 2x_5 & + & 6x_6 & \leq & 4 \\
    &  &  & x_2 & + & x_3 & - & 2x_4 & - & x_5 & + & 2x_6 & \leq & 5 \\
    & 4x_1 & - & 2x_2 & - & x_3 & + & 3x_4 & + & x_5 & - & x_6 & \leq & 3 \\
    & 3x_1 &  &  &  &  & + & x_4 &  &  & + & 2x_6 & \leq & 9 \\
    & \multicolumn{13}{l}{x_1,x_2,x_3,x_4,x_5,x_6\geq 0}
\end{array}
\]

\begin{enumerate}[a)]
    \item Plantear el problema dual.
    \item Utilizando el problema dual, mostrar que $(\frac{5}{3}, 0, \frac{7}{3}, 0, 0, \frac{4}{3})$ es solución óptima del problema primal.
    \item Hallar los valores de $\alpha,\beta\in \mathbb{R}$ para los cuales $(\frac{5}{3}, 0, \frac{7}{3}, 0, 0, \frac{4}{3})$ es solución óptima del problema con las mismas restricciones y con función objetivo: \[x_1  + (\alpha + \beta) x_2 - x_3 + (1+\beta)x_4 - x_5 + 6x_6\]
\end{enumerate}
\end{quest}

%%% TERMINA EL EJERCICIO 3 %%%

%%% EMPIEZA EL EJERCICIO 4 %%%


\begin{quest} 
Utilizar el algoritmo de Branch \& Bound para resolver el siguiente problema de programaci\'on lineal entera, detallando cada divisi\'on en subproblemas durante el algoritmo mediante el esquema de \'arbol visto en clase. Comenzar ramificando en $x_2$ y explicitar el m\'etodo utilizado para la resoluci\'on de cada subproblema (ejemplo: m\'etodo gr\'afico, c\'alculo de v\'ertices de cada regi\'on y posteriormente evaluar en cada uno de ellos, SIMPLEX, etc\'etera). 

\[\begin{array}{rrrrrrrrc}
\max & 2x_1 & + & 3x_2 &  &   \\
\text{s.a:}	 & 4x_1 & + & 5x_2 & \leq & 27  \\
			 & x_1 & + & 4x_2 & \geq & 4 \\
			 &-x_1 & +  & x_2 & \leq & 3 \\
			 &	4x_1    & +  & x_2 & \leq & 21 \\
			 &	x_1    &  & & \geq & 0 \\
			 & & & x_2 & \geq & 0 \\
			 & \multicolumn{3}{r}{x_1,x_2} & \in & \mathbb{Z} 
\end{array}
\]



\end{quest}

%%% TERMINA EL EJERCICIO 4 %%%

%%% TERMINA EL PARCIAL %%%



\end{document}

